% Fragmento integrable. Preambulo anfitrion: amsmath, amssymb, longtable, hyperref.
\section{Entrée transverse et changement d’échelle local}
\label{hmtfg:fragmento:entrada-transversal-y-escalado-local}

\subsection{Provenance et résultat}
\label{hmtfg:escalado:1}

APP évalue la somme et le produit sur les deux feuilles du réseau, en conservant quotient et reste. TRIT détermine le régime et l’orientation. Le transport TPK sélectionne, transporte et actualise l’état avec ses registres de retenue, de frontière et de mémoire. Ses lecteurs agissent sur la même structure discrète du continu. Le prolongement \texttt{w6→w12→w18→w24→w30→R36→G9} conserve cette provenance et distingue le retour de phase de l’incrément de mémoire.

Le lecteur critique de la roue dodécaphasique, construit dans les antécédents operatoriels et transporté par le raffinement compatible, produit
\[
\phi_j=\frac{(3j-1)\pi_{\mathrm{HMT}}}{18},\qquad
w_j=\frac1{12},\qquad
k_j=\sqrt{\frac2{\sum_{\ell=1}^{12}w_\ell\phi_\ell^2}}\,\phi_j
=\frac{2(3j-1)}{\sqrt{899}}.
\]
La coordonnée \(\pi_{\mathrm{HMT}}\) conserve son origine APP–TRIT–TPK ; son facteur commun s’annule dans ce deuxième moment. Le lecteur résultant est
\begin{equation}
u_0(x)=\frac1{12}\sum_{j=1}^{12}(1-\cos k_jx),\qquad
f_\lambda(x)=1-\lambda u_0(x).
\label{hmtfg:escalado:eq:1}
\end{equation}
Le secteur focal est le secteur uniforme, \(\eta=0\). L’extension antérieure conserve séparément la famille pondérée \(|\eta|\le1/40\). Les nouvelles bornes numériques correspondent à \eqref{hmtfg:escalado:eq:1}.

La famille est paire, entière, avec \(f_\lambda(0)=1\), un point critique quadratique et un deuxième moment \(\sum_jw_jk_j^2=2\). Le paramètre \(\lambda\) parcourt une famille postérieure au générateur HMT. La constante cible de Feigenbaum reste extérieure aux entrées de \eqref{hmtfg:escalado:eq:1}.

\textbf{Résultat principal.} Il existe un unique
\begin{equation}
\lambda_*\in[1.874038,1.874039]
\label{hmtfg:escalado:eq:2}
\end{equation}
pour lequel le quatrième retour normalisé de \eqref{hmtfg:escalado:eq:1} appartient à la feuille stable locale du point fixe réel de doublement. Le croisement de la famille avec cette feuille est transverse, avec une séparation quantifiée. Pour tout \(n\ge0\),
\begin{equation}
\boxed{\|R^{4+n}f_{\lambda_*}-g\|_{\mathcal A}
<0.001666\,(0.83996)^n.}
\label{hmtfg:escalado:eq:3}
\end{equation}
La cascade superstable locale associée à ce croisement possède des paramètres \(\mu_j\) et des constantes \(C\ne0\), \(0<\theta<1\), tels que
\begin{equation}
\mu_j-\lambda_*=C\delta_{\mathrm F}^{-j}\bigl[1+O(\theta^j)\bigr],
\qquad
\frac{\mu_{j-1}-\mu_{j-2}}{\mu_j-\mu_{j-1}}\longrightarrow\delta_{\mathrm F}.
\label{hmtfg:escalado:eq:4}
\end{equation}
Ici \(\delta_{\mathrm F}>1\) est la valeur propre instable de l’opérateur de doublement. La sélection de \(\mu_j\) s’effectue par les retours locaux de la section~\ref{hmtfg:escalado:8}. L’identification avec le sélecteur global est démontrée dans les sections~\ref{hmtfg:escalado:15}--\ref{hmtfg:escalado:17} : tous les termes globaux conservent la combinatoire et les centres coïncident à tout niveau suffisamment élevé lorsque leurs périodes sont égalisées.


\subsection{Carte analytique et retour effectif}
\label{hmtfg:escalado:2}

On utilise la carte de fonctions paires
\[
s=x^2,\qquad z=\frac{s-1}{5/2},\qquad
f(x)=1-x^2V(z),\qquad V(z)=\sum_{j\ge0}v_jz^j.
\]
Les coordonnées de l’espace de Banach sont
\begin{equation}
u=10v_0,\qquad \nu=(v_1,v_2,\ldots),\qquad
\|(\Delta u,\Delta\nu)\|_{\mathcal A}
=|\Delta u|+\sum_{j\ge1}|\Delta v_j|.
\label{hmtfg:escalado:eq:5}
\end{equation}
En particulier, le coefficient constant de \(V\) a un poids de dix. Cette normalisation coïncide avec la carte \(u/10\) de Hertling–Spandl, section 2 de la source.

Soit \(a=-f(1)=v_0-1\). Pour les retours réels admissibles,
\begin{equation}
(Rf)(x)=-\frac{f(f(-ax))}{a}.
\label{hmtfg:escalado:eq:6}
\end{equation}
Le programme évalue une factorisation exacte de \eqref{hmtfg:escalado:eq:6}. En définissant
\[
S=\frac{a^2s-1}{5/2},\qquad b=V(S),\qquad
h=1-a^2sb,\qquad t=\frac{h^2-1}{5/2},
\]
\[
\operatorname{shift}V(z)=\sum_{j\ge0}v_{j+1}z^j,
\]
on obtient
\begin{equation}
\boxed{V_R=a\,b\,(h+1)\left[V(t)+\frac25\operatorname{shift}V(t)\right].}
\label{hmtfg:escalado:eq:7}
\end{equation}
Pour la vérifier, on utilise \(v_0=1+a\),
\(h^2-1=-a^2sb(h+1)\) et
\(V(t)-v_0=t\operatorname{shift}V(t)\). Ainsi,
\(a+1-h^2V(t)=a^2sb(h+1)[V(t)+(2/5)\operatorname{shift}V(t)]\),
ce qui est précisément \(asV_R\).

Les moments initiaux sont rationnels :
\begin{equation}
\frac1{12}\sum_jk_j^{2m}
=\frac1{12}\left(\frac4{899}\right)^m
\sum_{j=1}^{12}(3j-1)^{2m}.
\label{hmtfg:escalado:eq:8}
\end{equation}
Ainsi, tant \eqref{hmtfg:escalado:eq:7} que la construction initiale utilisent des intervalles rationnels ; les fonctions trigonométriques initiales sont représentées par leurs séries avec une borne extérieure du reste.


\subsection{Bornes externes utilisées comme reconnaissance ultérieure}
\label{hmtfg:escalado:3}

Une fois \eqref{hmtfg:escalado:eq:1} construite, le théorème d’hyperbolicité de Lanford est utilisé comme certificat analytique de l’opérateur \eqref{hmtfg:escalado:eq:6}. Soient \(g_0\) le polynôme central publié et \(g\) le point fixe. Les estimations utilisées sont
\[
\|g-g_0\|_{\mathcal A}<\varepsilon_0=4\,10^{-5},
\]
et, dans la boule \(\|f-g_0\|_{\mathcal A}<0.01\), les blocs de \(DR\), notés \(R=(U,V)\), satisfont
\begin{equation}
U_u\ge a_0=4.521,\quad
\|U_\nu\|\le b_0=0.560,\quad
\|V_u\|\le c_0=0.756,\quad
\|V_\nu\|\le d_0=0.719.
\label{hmtfg:escalado:eq:9}
\end{equation}
Le point fixe possède une unique direction instable, de valeur propre réelle positive \(\delta_{\mathrm F}>1\) ; le reste du spectre se situe à l’intérieur du disque unité. Les chiffres de \eqref{hmtfg:escalado:eq:9} proviennent des estimations publiées, et non d’une nouvelle reproduction de la preuve de Lanford dans ce dossier. Les coefficients de \(g_0\) sont explicitement incorporés au polynôme central de référence du programme, conformément au tableau 2 de Hertling–Spandl.

Sources : \href{https://repo-archives.ihes.fr/A_GARDER/20180409/Test/I_Prepublications/LANFORD/1968-2013/P_81_17/P_81_17_web.pdf}{Lanford, estimations d’hyperbolicité}; \href{https://arxiv.org/pdf/1410.3277}{Hertling–Spandl, carte analytique et tableau 2}.


\subsection{Certificat rationnel d’entrée et de direction}
\label{hmtfg:escalado:4}

Le certificat d’entrée dans la renormalisation divise \eqref{hmtfg:escalado:eq:2} en seize intervalles rationnels égaux. Il propage quatre retours de \eqref{hmtfg:escalado:eq:7} et la dérivée exacte par rapport à \(\lambda\), en utilisant la différentiation de la composition et l’arrondi extérieur à cent décimales.

Chaque série se compose d’un polynôme de degré 32 à coefficients intervalles et d’une erreur analytique arbitraire \(E\), \(\|E\|_1\le\epsilon\). Cette erreur peut également modifier les coefficients de bas degré. Pour \(\|q\|_1<1\),
\begin{equation}
\|E\circ q\|_1\le\epsilon,\qquad
\|E'\circ q\|_1\le\frac{\epsilon}{(1-\|q\|_1)^2},\qquad
\|\operatorname{shift}E\|_1\le\epsilon.
\label{hmtfg:escalado:eq:10}
\end{equation}
Les produits comprennent les termes polynôme–erreur et erreur–erreur. La queue initiale est majorée par son premier terme en valeur absolue et une raison géométrique, en utilisant \(\|1+(5/2)z\|_1=7/2\). Chaque argument composé reste strictement à l’intérieur de la boule unité de cette algèbre.

Pour \(\Gamma(\lambda)=R^4f_\lambda=(u_\lambda,\nu_\lambda)\), le certificat donne les bornes conservatrices suivantes, uniformes en \eqref{hmtfg:escalado:eq:2} :

\begingroup\small
\begin{longtable}{p{0.465\linewidth}p{0.465\linewidth}}
\hline
Quantité & Borne extérieure certifiée \\
\hline
\(\|\Gamma(\lambda)-g_0\|_{\mathcal A}\) & \(<0.004993128\) \\
\(\|\nu_\lambda-\nu_{g_0}\|_1\) & \(<0.001396077\) \\
\(u'_\lambda\) & \(>3821.289115\) \\
\(\|\nu'_\lambda\|_1\) & \(<551.757000\) \\
\(\|\nu'_\lambda\|_1/u'_\lambda\) & \(<0.144386<1/4\) \\
Erreur analytique de la fonction & \(<8.482\,10^{-19}\) \\
Erreur analytique de la tangente & \(<4.569\,10^{-15}\) \\
\hline\end{longtable}
\endgroup

Les décimales complètes, les seize sous-intervalles et chaque retour intermédiaire sont conservés dans le certificat rationnel d’entrée.

Le cône \(\|\dot\nu\|_1\le|\dot u|/4\) est invariant par \eqref{hmtfg:escalado:eq:9} :
\[
|\dot u_{\rm nuevo}|\ge(4.521-0.560/4)|\dot u|
=4.381|\dot u|,
\]
\begin{equation}
\|\dot\nu_{\rm nueva}\|_1\le(0.756+0.719/4)|\dot u|
=0.93575|\dot u|<\tfrac14(4.381)|\dot u|.
\label{hmtfg:escalado:eq:11}
\end{equation}
La direction paramétrique certifiée entre donc dans un cône uniformément expansif de l’opérateur.


\subsection{Construction quantitative de la feuille stable}
\label{hmtfg:escalado:5}

Translatons \(g\) à l’origine et écrivons les coordonnées sous la forme \((x,y)=(u-u_g,\nu-\nu_g)\). Fixons
\begin{equation}
L=0.16,\qquad r=0.004,\qquad
A=a_0-Lc_0=4.40004,\qquad q=d_0+c_0L=0.83996.
\label{hmtfg:escalado:eq:12}
\end{equation}
Le cylindre \(|x|\le Lr\), \(\|y\|_1\le r\) se trouve à l’intérieur de la boule de \eqref{hmtfg:escalado:eq:9}, puisque
\((1+L)r+\varepsilon_0=0.00468<0.01\).

Considérons l’espace complet des fonctions \(\sigma:B_r\to\mathbb R\) avec \(\sigma(0)=0\) et de constante de Lipschitz au plus \(L\), muni de la distance uniforme. Pour chaque \(y\), on définit \((\mathcal G\sigma)(y)\) comme la racine dans \([-L\|y\|_1,L\|y\|_1]\) de
\begin{equation}
F_\sigma(x,y)=U(x,y)-\sigma(V(x,y)).
\label{hmtfg:escalado:eq:13}
\end{equation}
Sur cet intervalle, \(\|V(x,y)\|_1\le q\|y\|_1<r\). Lorsque \(x\) augmente, \eqref{hmtfg:escalado:eq:9} montre que \(F_\sigma\) croît avec une pente au moins égale à \(A\). À ses extrémités,
\[
F_\sigma(L\|y\|_1,y)
\ge[AL-b_0-Ld_0]\|y\|_1,
\]
et l’inégalité opposée vaut à l’extrémité négative. La marge est
\begin{equation}
AL-b_0-Ld_0=0.0289664>0.
\label{hmtfg:escalado:eq:14}
\end{equation}
On obtient une racine unique. En comparant \eqref{hmtfg:escalado:eq:13} pour deux valeurs de \(y\),
\[
\operatorname{Lip}(\mathcal G\sigma)
\le\frac{b_0+Ld_0}{A}
=\frac{0.67504}{4.40004}<0.16.
\]
En comparant deux graphes,
\begin{equation}
\|\mathcal G\sigma-\mathcal G\widetilde\sigma\|_\infty
\le A^{-1}\|\sigma-\widetilde\sigma\|_\infty,
\qquad A^{-1}<0.228.
\label{hmtfg:escalado:eq:15}
\end{equation}
Le théorème de contraction produit un unique graphe fixe \(x=\sigma_*(y)\). Sur celui-ci,
\begin{equation}
\|y_n\|_1\le q^n\|y_0\|_1,
\qquad
\|(x_n,y_n)\|_{\mathcal A}\le(1+L)q^n\|y_0\|_1.
\label{hmtfg:escalado:eq:16}
\end{equation}
De plus, la distance verticale \(d(x,y)=x-\sigma_*(y)\) satisfait
\[
|d(R(x,y))|\ge A|d(x,y)|
\]
tant que l’orbite reste dans le cylindre. Toute orbite qui y demeure indéfiniment appartient au graphe. Celui-ci identifie exactement la feuille stable locale et fournit l’estimation \eqref{hmtfg:escalado:eq:16}. Sa régularité locale coïncide avec celle de la variété stable analytique du point fixe hyperbolique.


\subsection{Existence, unicité et transversalité du paramètre}
\label{hmtfg:escalado:6}

Pour \(\lambda\) dans \eqref{hmtfg:escalado:eq:2}, définissons
\[
H(\lambda)=u_\lambda-u_g-
\sigma_*(\nu_\lambda-\nu_g).
\]
Le certificat assure \(\|\nu_\lambda-\nu_g\|_1<0.001436077<r\). Pour \(\lambda_2>\lambda_1\),
\begin{equation}
H(\lambda_2)-H(\lambda_1)
\ge\int_{\lambda_1}^{\lambda_2}
\bigl[u'_t-L\|\nu'_t\|_1\bigr]dt
>3733.007995\,(\lambda_2-\lambda_1).
\label{hmtfg:escalado:eq:17}
\end{equation}
Pour déterminer les signes aux extrémités, on utilise \(\|g-g_0\|_{\mathcal A}<\varepsilon_0\) et \(|\sigma_*(y)|\le L\|y\|_1\). L’évaluation ponctuelle rationnelle produit
\begin{equation}
H(1.874038)<-0.0011856110945,
\qquad H(1.874039)>0.0021866053399.
\label{hmtfg:escalado:eq:18}
\end{equation}
Par continuité, il existe une racine ; \eqref{hmtfg:escalado:eq:17} prouve son unicité et la transversalité du croisement. À cette racine, \eqref{hmtfg:escalado:eq:16} et
\((1+L)(0.001396077+0.00004)<0.001666\)
démontrent \eqref{hmtfg:escalado:eq:3}.

La borne est uniforme en le nombre de retours : elle convertit l’avancement en profondeur en une erreur operatorielle géométrique explicite.


\subsection{Admissibilité réelle de tous les retours}
\label{hmtfg:escalado:7}

Les quatre retours initiaux s’accompagnent, sur chaque sous-intervalle rationnel, des inégalités
\begin{equation}
0<a=-f(1)<1,\qquad b=f(a)>a,\qquad f(b)<a.
\label{hmtfg:escalado:eq:19}
\end{equation}
Le programme conserve \(a,b,f(b)\) avant chaque application de \eqref{hmtfg:escalado:eq:7}. L’unimodalité, la parité et le point critique quadratique sont conservés par ces retours normalisés.

Pour les itérations suivantes, une vérification uniforme dans la boule de \eqref{hmtfg:escalado:eq:9} suffit. Si \(E=V-V_{g_0}\), la norme \eqref{hmtfg:escalado:eq:5} donne, pour \(-0.4\le z\le0\),
\[
|E(z)|\le0.004,\qquad |E'(z)|\le0.01.
\]
La seconde inégalité utilise \(\sup_{j\ge1}j(0.4)^{j-1}=1\). L’évaluation rationnelle des coefficients centraux fournit
\begin{equation}
\frac65<V(z)<\frac85,\qquad |V'(z)|<\frac12,
\qquad 0.39<a<0.41.
\label{hmtfg:escalado:eq:20}
\end{equation}
En conséquence,
\[
f'(x)=-2x\left[V(z)+\frac25x^2V'(z)\right],
\]
et l’expression entre crochets dépasse un pour \(0\le x\le1\). L’application est unimodale et préserve \([-1,1]\). De plus,
\[
b=f(a)>1-(0.41)^2\frac85=\frac{4569}{6250}>a,
\]
\begin{equation}
f(b)<1-\left(\frac{4569}{6250}\right)^2\frac65
=\frac{35028967}{97656250}<0.39<a.
\label{hmtfg:escalado:eq:21}
\end{equation}
Les orbites de la feuille stable satisfont \eqref{hmtfg:escalado:eq:19} à toute profondeur. Ainsi, les itérations analytiques de l’opérateur correspondent à des retours réels de doublement, dont les domaines et les orientations sont conservés.


\subsection{Changement d’échelle de la cascade superstable locale}
\label{hmtfg:escalado:8}

La conclusion métrique exige deux croisements : celui de la famille avec la feuille stable, démontré dans la section~\ref{hmtfg:escalado:6}, et celui de la variété instable avec la cible superstable. Le second est démontré par Eckmann–Wittwer pour
\(\Sigma_2=\{\psi:\psi(1)=0\}\), correspondant au cycle superstable de période deux. La normalisation \eqref{hmtfg:escalado:eq:6} coïncide avec celle de l’opérateur réel pair utilisé dans ce travail.

Le théorème général de Collet–Eckmann–Lanford, section 6 de la source citée, proposition 6.1 et théorèmes 6.2–6.3, s’applique à une application d’un espace de Banach \(C^2\), un point fixe possédant une unique valeur propre instable \(\delta_{\mathrm F}>1\), une courbe transverse à sa feuille stable et une cible transverse à sa variété instable. Les préimages locales de la cible rencontrent la courbe en des points qui satisfont
\begin{equation}
\delta_{\mathrm F}^j(\mu_j-\lambda_*)\longrightarrow C\ne0.
\label{hmtfg:escalado:eq:22}
\end{equation}
L’indice \(j\) compte les retours locaux. Comme la courbe initiale est ici \(R^4f_\lambda\), la période physique est \(2^{j+5}\) lorsque la cible est prise dans \(\Sigma_2\) ; tout nombre fixe de pas utilisé pour ramener la cible dans le voisinage local est incorporé à un décalage fixe de l’indice et à \(C\).

Les conditions réelles de la section~\ref{hmtfg:escalado:7} et la branche réelle de doublement de la cible fixent la période exacte et l’itinéraire. La suite est sélectionnée par ce prolongement local orienté. En prenant les différences dans \eqref{hmtfg:escalado:eq:22}, on obtient la limite des rapports de \eqref{hmtfg:escalado:eq:4}.

Sources des deux résultats utilisés : \href{https://link.springer.com/article/10.1007/BF01013368}{Eckmann–Wittwer, \emph{A complete proof of the Feigenbaum conjectures}}; \href{https://people.math.harvard.edu/~knill/history/lanford/papers/ColletEckmannLanford.pdf}{Collet–Eckmann–Lanford, section 6}.

\subsubsection{Reste asymptotique en puissance}

L’analyticité de notre famille permet de préciser \eqref{hmtfg:escalado:eq:22}. Dans la construction de coordonnées normales de Collet–Eckmann–Lanford, seules la variété stable, la variété instable et l’hypersurface cible sont aplaties. Ces trois sous-variétés admettent des changements locaux \(C^2\) à dérivées bornées. La coupure lisse employée s’effectue dans la coordonnée scalaire instable ; le feuilletage complet s’obtient ensuite par le changement construit.

Dans ces coordonnées, la correction s’écrit comme une série d’incréments \(w_j\) dont la norme \(C^1\) satisfait
\[
\|w_j\|_{C^1}\le C_1\rho^j,\qquad 0<\rho<1.
\]
La règle de dérivation en chaîne pour l’application munie d’une coupure, avec des dérivées d’ordres un et deux uniformément bornées, fournit
\[
\|w_j\|_{C^2}\le C_2B^j
\]
pour un certain \(B\ge1\). La norme pondérée de la construction contrôle la norme \(C^1\) ordinaire sur le domaine considéré. Par interpolation,
\[
[Dw_j]_\beta\le C_3
\left(\rho^{1-\beta}B^\beta\right)^j.
\]
On choisit
\begin{equation}
0<\beta<\min\left\{1,
\frac{-\log\rho}{\log B-\log\rho}\right\}.
\label{hmtfg:escalado:eq:23}
\end{equation}
La série converge dans \(C^{1,\beta}\). Soit \(z\) la coordonnée normale obtenue, normalisée de sorte que les préimages de la cible satisfassent \(z=\delta_{\mathrm F}^{-j}\), en absorbant le décalage fixe de l’indice. Comme notre courbe est analytique et transverse,
\[
z(\Gamma(\lambda))=c(\lambda-\lambda_*)
+O(|\lambda-\lambda_*|^{1+\beta}),\qquad c\ne0.
\]
L’inversion locale donne
\begin{equation}
\mu_j-\lambda_*=C\delta_{\mathrm F}^{-j}
+O\left(\delta_{\mathrm F}^{-(1+\beta)j}\right).
\label{hmtfg:escalado:eq:24}
\end{equation}
Cela démontre le reste de \eqref{hmtfg:escalado:eq:4}, avec \(\theta=\delta_{\mathrm F}^{-\beta}\). L’énoncé affirme l’existence des constantes du reste paramétrique. Le taux operatoriel explicite \(0.83996\) de \eqref{hmtfg:escalado:eq:3} contrôle, quant à lui, la convergence des applications renormalisées.

\subsubsection{Conversion du reste en erreur sur les rapports}

Si \(\mu_j-\lambda_*=C\delta_{\mathrm F}^{-j}(1+e_j)\), avec \(|e_j|\le M\theta^j\), on écrit
\[
\mu_{j-1}-\mu_j=C(\delta_{\mathrm F}-1)\delta_{\mathrm F}^{-j}(1+r_j),
\quad
r_j=\frac{\delta_{\mathrm F} e_{j-1}-e_j}{\delta_{\mathrm F}-1}.
\]
Avec \(K=M(\delta_{\mathrm F}+\theta)/(\delta_{\mathrm F}-1)\), on a \(|r_j|\le K\theta^{j-1}\). Par conséquent,
\begin{equation}
\boxed{
\left|\frac{\mu_{j-2}-\mu_{j-1}}{\mu_{j-1}-\mu_j}-\delta_{\mathrm F}\right|
\le\frac{\delta_{\mathrm F} K(1+\theta)\theta^{j-2}}
{1-K\theta^{j-1}}}
\label{hmtfg:escalado:eq:25}
\end{equation}
lorsque \(K\theta^{j-1}<1\). Cette formule sépare la précision d’évaluation de chaque racine de l’erreur de troncature de la cascade.


\subsection{Prolongement certifié des huit racines initiales}
\label{hmtfg:escalado:9}

Le certificat de prolongement fini conserve la famille \eqref{hmtfg:escalado:eq:1} et le certificat antérieur de huit racines. Pour \(2\le n\le8\), il étudie
\[
S_n(\lambda)=f_\lambda^{2^n}(0)
\]
depuis la boîte certifiée de \(\lambda_{n-1}\) jusqu’à \(1.874039\). Le recouvrement rationnel complet contient exactement la racine héritée \(\lambda_{n-1}\) et la nouvelle racine \(\lambda_n\). En particulier, \(\lambda_n\) est l’unique nouvelle racine dans \((\lambda_{n-1},1.874039]\).

Les boîtes sont exclues par le signe du retour ou par une dérivée locale de signe certifié. Les boîtes d’ancrage conservent l’existence et l’unicité par les signes aux extrémités et la monotonie locale. Leurs extrémités couvrent exactement chaque intervalle ; les adjacences sont vérifiées. L’évaluation utilise les moments \eqref{hmtfg:escalado:eq:8}, l’évaluation de Horner d’ordre 32, des queues géométriques extérieures et \(|u_0''|\le2\). La queue uniforme du potentiel est inférieure à \(1.294\,10^{-58}\) sur \(|x|\le3/2\).

Le certificat final contient 650 boîtes traitées et 332 feuilles de recouvrement. Le premier niveau possède l’expression exacte \(\lambda_1=1/u_0(1)\). Les approximations suivantes sont présentées uniquement à titre indicatif ; leurs boîtes complètes sont conservées dans les certificats :

\begingroup\small
\begin{longtable}{p{0.465\linewidth}p{0.465\linewidth}}
\hline
Niveau & Paramètre superstable \\
\hline
1 & 1.345913990471832792210949095… \\
2 & 1.763379787846910127290794030… \\
3 & 1.850413796950136464530566778… \\
4 & 1.868979756606798125150574309… \\
5 & 1.872955034435659740004205128… \\
6 & 1.873806367467030119412262311… \\
7 & 1.873988695221365362307168780… \\
8 & 1.874027744154450672897007573… \\
\hline\end{longtable}
\endgroup
